% ========================================================
% BÀI TẬP TOÁN 5 - KẾT NỐI TRI THỨC VỚI CUỘC SỐNG
% Bài 69: Ôn tập các phép tính với số tự nhiên, phân số, số thập phân
% Định dạng: BTTeX Part III - Trắc nghiệm điền số ngắn
% ========================================================

% Câu 1
\begin{ex}
  Tính: $1\,234 + 5\,678$. Kết quả là bao nhiêu?
  \shortans{6912}
  \loigiai{
    $1\,234 + 5\,678 = 6\,912$.
  }
\end{ex}

% Câu 2
\begin{ex}
  Tính: $10\,000 - 3\,456$. Kết quả là bao nhiêu?
  \shortans{6544}
  \loigiai{
    $10\,000 - 3\,456 = 6\,544$.
  }
\end{ex}

% Câu 3
\begin{ex}
  Tính: $125 \times 8$. Kết quả là bao nhiêu?
  \shortans{1000}
  \loigiai{
    $125 \times 8 = 1\,000$.
  }
\end{ex}

% Câu 4
\begin{ex}
  Tính: $4\,500 : 90$. Kết quả là bao nhiêu?
  \shortans{50}
  \loigiai{
    $4\,500 : 90 = 50$.
  }
\end{ex}

% Câu 5
\begin{ex}
  Tính: $\frac{2}{5} + \frac{3}{5}$. Kết quả là bao nhiêu?
  \shortans{1}
  \loigiai{
    $\frac{2+3}{5} = \frac{5}{5} = 1$.
  }
\end{ex}

% Câu 6
\begin{ex}
  Tính: $\frac{7}{10} - \frac{1}{5}$. Kết quả viết dưới dạng phân số tối giản có mẫu số là bao nhiêu?
  \shortans{2}
  \loigiai{
    $\frac{7}{10} - \frac{2}{10} = \frac{5}{10} = \frac{1}{2}$. Mẫu số là 2.
  }
\end{ex}

% Câu 7
\begin{ex}
  Tính: $\frac{3}{4} \times \frac{8}{9}$. Kết quả viết dưới dạng phân số tối giản có mẫu số là bao nhiêu?
  \shortans{3}
  \loigiai{
    $\frac{3 \times 8}{4 \times 9} = \frac{2}{3}$. Mẫu số là 3.
  }
\end{ex}

% Câu 8
\begin{ex}
  Ở câu trên, tử số là bao nhiêu?
  \shortans{2}
  \loigiai{
    Tử số là 2.
  }
\end{ex}

% Câu 9
\begin{ex}
  Tính: $\frac{5}{6} : \frac{5}{12}$. Kết quả là bao nhiêu?
  \shortans{2}
  \loigiai{
    $\frac{5}{6} \times \frac{12}{5} = 2$.
  }
\end{ex}

% Câu 10
\begin{ex}
  Tính: $12{,}5 + 7{,}5$. Kết quả là bao nhiêu?
  \shortans{20}
  \loigiai{
    $12{,}5 + 7{,}5 = 20$.
  }
\end{ex}

% Câu 11
\begin{ex}
  Tính: $50 - 18{,}4$. Kết quả nhân 10 được bao nhiêu?
  \shortans{316}
  \loigiai{
    $50 - 18{,}4 = 31{,}6$. Nhân 10 được 316.
  }
\end{ex}

% Câu 12
\begin{ex}
  Tính: $2{,}5 \times 4$. Kết quả là bao nhiêu?
  \shortans{10}
  \loigiai{
    $2{,}5 \times 4 = 10$.
  }
\end{ex}

% Câu 13
\begin{ex}
  Tính: $15 : 0{,}5$. Kết quả là bao nhiêu?
  \shortans{30}
  \loigiai{
    $15 : 0{,}5 = 30$.
  }
\end{ex}

% Câu 14
\begin{ex}
  Tính nhanh: $25 \times 36 \times 4$. Kết quả là bao nhiêu?
  \shortans{3600}
  \loigiai{
    $(25 \times 4) \times 36 = 100 \times 36 = 3\,600$.
  }
\end{ex}

% Câu 15
\begin{ex}
  Tính nhanh: $4{,}8 \times 15 + 4{,}8 \times 85$. Kết quả là bao nhiêu?
  \shortans{480}
  \loigiai{
    $4{,}8 \times (15 + 85) = 4{,}8 \times 100 = 480$.
  }
\end{ex}

% Câu 16
\begin{ex}
  Tìm $x$: $x + 45 = 120$. Giá trị của $x$ là bao nhiêu?
  \shortans{75}
  \loigiai{
    $x = 120 - 45 = 75$.
  }
\end{ex}

% Câu 17
\begin{ex}
  Tìm $x$: $x - 3{,}5 = 6{,}5$. Giá trị của $x$ là bao nhiêu?
  \shortans{10}
  \loigiai{
    $x = 6{,}5 + 3{,}5 = 10$.
  }
\end{ex}

% Câu 18
\begin{ex}
  Tìm $x$: $x \times 8 = 720$. Giá trị của $x$ là bao nhiêu?
  \shortans{90}
  \loigiai{
    $x = 720 : 8 = 90$.
  }
\end{ex}

% Câu 19
\begin{ex}
  Tìm $x$: $450 : x = 15$. Giá trị của $x$ là bao nhiêu?
  \shortans{30}
  \loigiai{
    $x = 450 : 15 = 30$.
  }
\end{ex}

% Câu 20
\begin{ex}
  Tính: $3{,}6 \times 100$. Kết quả là bao nhiêu?
  \shortans{360}
  \loigiai{
    $3{,}6 \times 100 = 360$.
  }
\end{ex}

% Câu 21
\begin{ex}
  Tính: $45 : 10$. Kết quả nhân 10 là bao nhiêu?
  \shortans{45}
  \loigiai{
    $45 : 10 = 4{,}5$. Nhân 10 được 45.
  }
\end{ex}

% Câu 22
\begin{ex}
  Tính giá trị biểu thức: $120 - (30 + 20 \times 2)$. Kết quả là bao nhiêu?
  \shortans{50}
  \loigiai{
    $120 - (30 + 40) = 120 - 70 = 50$.
  }
\end{ex}

% Câu 23
\begin{ex}
  Tính: $(1{,}2 + 2{,}8) \times 5$. Kết quả là bao nhiêu?
  \shortans{20}
  \loigiai{
    $4 \times 5 = 20$.
  }
\end{ex}

% Câu 24
\begin{ex}
  Một tổ may 12 bộ quần áo hết 36 mét vải. May 20 bộ như thế hết bao nhiêu mét vải?
  \shortans{60}
  \loigiai{
    Mỗi bộ hết $36 : 12 = 3\text{ m}$. $20$ bộ hết $3 \times 20 = 60\text{ m}$.
  }
\end{ex}

% Câu 25
\begin{ex}
  Mua 5 quyển vở hết 45 nghìn đồng. Mua 10 quyển vở như thế hết bao nhiêu nghìn đồng?
  \shortans{90}
  \loigiai{
    $45 \times 2 = 90$ nghìn đồng.
  }
\end{ex}
